\documentclass{article}
\usepackage{amsmath,amssymb,hyperref,url}
\title{Do contracts discourage permitted trades?\\An unfinished connexion between markets and cooperation}
\author{}
\date{}
\begin{document}
\maketitle
\begin{abstract}
This account connects \emph{Competitive Market Behavior of LLMs} (Q) with \emph{Commitment To Cooperation With
Self-Negotiated Contracts} (P). The peers asked whether contracts could improve language-model trading, then
narrowed the question to whether contractual wording makes agents infer restrictions that the rules do not
impose. They established limits on payment-based explanations, feasible cooperation in P, and the distinction
between market surplus and price convergence. They also corrected an overstrong reward bound: payments to a
nonfinisher may enlarge the set of outcomes that benefit both players. The thread ended at PAUSE because these
results did not establish the proposed shared cause, and the controlled agent trials needed to test it had not
been run.
\end{abstract}

\section{Introduction}
Q studies language models acting as buyers and sellers in a double auction, where either side can post an
offer \cite{Q}. Each agent can trade once per round. A buyer who offers at least the current asking price
trades immediately at that price; the seller's case is symmetric. Both parties then leave the round. Q reports
inefficient allocations and small price adjustments despite profitable trading opportunities. It asks whether
a familiar market institution works well when language models replace human traders.

P studies bargaining and navigation in Colored Trails \cite{P}. Players spend coloured chips to move towards a
goal and can help each other by exchanging chips or covering moves. Regular trading exchanges chips directly.
In the alternative Pay-for-Partner setting, a promise to cover a later move can be refused when payment falls
due. P studies contracts that automate tile payments, natural-language agreements interpreted during play, and
point transfers conditional on finishing. Contracts can improve cooperation, but some conditions also suppress
further trade.

The initial connexion was to bring P's contracts into Q. Ada and Emmy rejected a simple explanation based on
making settlement credible: Q already executes accepted trades automatically. They instead asked whether
contractual wording could make an agent wrongly treat an existing obligation as a restriction on other
profitable actions. P supplies agreements that leave further coordination necessary. Q supplies an already
enforced market where wording could change without adding protection against later refusal. This is a proposed
comparison, not a demonstrated common mechanism.

\section{What was done}
The peers first compared rules and reported behaviour. On P's mutually dependent boards under regular trading,
natural-language contracts give normalised joint reward of \(0.64\), against \(0.93\) without contracts.
Normalisation divides the combined score by P's benchmark of \(140\). Compared with programmatic tile
contracts, natural-language contracts have similar mean executed coverage, about \(2.6\) moves, but fewer
trades: \(1.95\) against \(3.28\). These observations suggested that contract use could discourage useful
activity. The peers kept them separate from delayed Pay-for-Partner results, where refusal to honour a promise
is a different problem.

They next checked what contracts could change economically. Ada reduced a restricted bilateral payment in Q to
an effective price. Emmy constructed chip exchanges that let both players finish in P, and examples separating
efficient allocation from price convergence in Q. Ada checked the constructions independently. These
calculations distinguished failure to find a feasible gain from the absence of such a gain.

The peers then designed a comparison that would hold executable terms fixed and vary their display. In Q, an
unfilled order would appear either as an order record or as a conditional obligation. A disclosed final
trading opportunity would remove sensible waiting as an explanation for refusal. Separate questions would test
which actions the agent believed were permitted, and whether a permission clarification corrected its error.

Source checks narrowed possible novelty. Hantel's author summary already describes price anchors, market
failures, and placebo and debiasing controls \cite{hantel}. Ada inspected the summary rather than the full
linked paper; the associated SSRN URL remained a search lead \cite{ssrn}. Sajja and colleagues'
\emph{Evaluating Rational Contracting in Natural Language}, or ContractSim, already studies contracts as
policy constraints and evaluates rational performance \cite{contractsim}. These checks did not establish that
the narrower experiment was new.

Finally, Emmy reopened the reward calculation. The earlier claim that \(140\) was an unconditional maximum in
P was too strong. Concentrating chips in a single finisher could exceed it, and a points contract might
compensate the other player for not finishing. Both peers incorporated the correction. Their attempted
implementation audit failed, so the revised claim remained conditional on the written rules.

\section{Results and their support}
\subsection{A same-trade payment can be an effective price}
Let \(v\) be a buyer's value, \(c\) a seller's cost, \(p\) the transaction price, and \(d\) a net payment from
that buyer to that seller. Assume that \(d\) is triggered only by their same atomic trade, with no financing
constraint or additional obligation. With quasilinear utility, meaning money enters payoffs linearly, their
payoffs \(u_B\) and \(u_S\) are
\[
u_B=v-p-d,\qquad u_S=p+d-c.
\]
The effective price \(p_{\mathrm{eff}}=p+d\) reproduces the allocation and realised payoffs. Combined surplus
remains \(u_B+u_S=v-c\).

Ada gave this reduction in \path{ada/research-note.md}; Emmy accepted its qualified form in
\path{emmy/boundaries.md}. It does not establish equivalence of strategies. Payments tied to particular
partners can reorder offers, commitments can change information or available actions, and payments triggered
by third-party trades create different incentives.

\subsection{Ordinary exchange reaches the both-finish bound}
Each player in P begins with fourteen own-colour chips and two green chips. A six-move shortest route consumes
five red or blue chips and a green chip for the goal. Exchanging five red chips for five blue chips gives each
player nine own-colour chips, five opposite-colour chips and two green chips. Both can afford every shortest
route. Each finishes with ten chips and scores \(20+5(10)=70\).

Let \(R_{\mathrm{joint}}\) denote the sum of the final scores. If both players finish, at least twelve of the
initial thirty-two chips are consumed. The goal awards total forty points and each remaining chip is worth
five, so
\[
R_{\mathrm{joint}}\leq 2(20)+5(32-12)=140.
\]
Without point transfers, a nonfinisher receives zero. Since both no-interaction baselines are nonnegative,
strict improvement over both baselines requires both players to finish. The same bound therefore applies to
strict mutual improvement without point transfers.

On an asymmetric board, take Red to be the player who can finish alone and Blue the player who cannot.
Exchanging five red chips for six blue chips leaves Red with nine red, six blue and two green chips, and Blue
with five red, eight blue and two green chips. Both can still take any shortest route. Red scores \(75\) and
Blue scores \(65\), strictly exceeding their respective no-interaction baselines of at most \(70\) and zero.

Emmy's \path{emmy/verify_bounds.py} checked all \(3432\) balanced boards and \(20\) shortest paths per board
against five inventories, including the concentration case below. Its \(343200\) checks support resource
feasibility. Ada independently checked the swap bounds. These are feasible constructions, not claims that
agents choose them or that bargaining leads to them in equilibrium.

\subsection{Compensating a nonfinisher changes the bound}
Under P's written arbitrary-quantity exchange rule, Red can give one red chip for Blue's fourteen blue and two
green chips. Red then holds thirteen red, fourteen blue and four green chips; Blue holds only one red chip.
Red can finish with twenty-five chips left and score \(145\). Blue fails and scores zero. This refutes an
unconditional \(140\) bound under those rules, though it does not benefit Blue.

Suppose a programmatic points contract pays a recipient who does not finish. Red can first promise Blue ten
points upon Red's completion. The same exchange then yields \((135,10)\), with Red's score first and Blue's
second. Both strictly beat their asymmetric no-interaction baselines and joint reward is \(145\). This does
not improve both scores relative to the ordinary exchange outcome \((75,65)\): Blue receives less.

Emmy proposed this correction and Ada confirmed its arithmetic and textual basis. Point contracts may
therefore expand the set of strictly baseline-improving rewards by compensating noncompletion. The earlier
claim that P's gains must concern behaviour alone was withdrawn. Joint reward, individual improvement and the
number of finishers must be reported separately.

The implementation qualification matters. P's prompt says a nonfinisher receives zero total points, which
needs reconciliation with its point-transfer rule. Chip quantities, scoring order and termination also need
checking. Ada could not open P's anonymous repository \cite{prepo}; its API returned \texttt{not\_connected}
\cite{papi}. The resource check does not establish how the code handles this construction or demonstrate a
repository defect.

\subsection{Market surplus and price convergence differ}
In Q, surplus means buyer value minus seller cost, summed over trades. The ranked gains are
\(2.5,2,1.5,1,0.5,0\), so maximum surplus is \(7.5\). Five positive-gain trades suffice. Emmy paired the five
highest-value buyers with the five lowest-cost sellers at prices equal to the sellers' costs:
\(0.75,1,1.25,1.5,1.75\). This gives full surplus despite a price-dispersion coefficient of about \(41.46\)
under Q's measure of distance from its competitive price of \(2\).

Conversely, one trade between a buyer with value \(3.25\) and a seller with cost \(0.75\), at price \(2\), has
zero dispersion but only one third of maximum surplus. Emmy checked these examples in the same script; Ada
checked maximum surplus independently. They show why price convergence cannot substitute for surplus
measurement. They do not dispute Q's measured efficiency losses or describe typical play.

\section{Objections and unresolved explanations}
Equal mean covered moves in P does not imply equal contract quality, routes, inventories, timing or remaining
gains. Executed coverage is itself affected by treatment. P's audit examined approved moves, so it cannot
bound mistaken rejections. Holding terms fixed and using identical deterministic enforcement were proposed
remedies, but were not carried out.

Q's single-trade rule also limits the comparison. An agent that has traded leaves the round and cannot
undertake P-like further trading after a partial agreement. The proposed Q test therefore concerns unfilled
orders. Moreover, a profitable trade need not be the best choice if waiting could bring a better price. Only
the disclosed final-opportunity probe removes that explanation; this disclosure would be an addition to Q's
experiment.

Most directly, saying further trade is unnecessary does not mean believing it is forbidden. P's explanations
could reflect a mistaken belief that the existing plan is sufficient. The peers proposed separate tests of
permission and plan sufficiency. Attention, planning, context burden and interpretation remained alternatives.
Changed trading rates alone would establish sensitivity to presentation, not the proposed mistaken
restriction.

Practical questions remained too. Q's repository README says that the configuration seed controls
zero-intelligence trader randomness while mechanism scheduling is unseeded \cite{qrepo}. Ada warned that
matching seeds would not produce paired market histories. Later trials would need explicit scheduling coupling
or independently randomised market runs. The peers had neither effect sizes nor a justified sample size, and
ran no language-model experiment or market replication.

\section{Why the thread stopped}
The verifier recorded PAUSE in its final and only supplied round. It accepted that the note incorporated the
corrections, but judged that the economic boundaries did not establish a non-obvious result about the
connexion. P's unnecessary-trade judgments did not demonstrate false prohibitions; Q's incremental offers did
not identify contract-induced restrictions. The record contained a design rather than the controlled agent
evidence needed to support it, and the access needed to obtain that evidence was unestablished.

The smallest restart would audit a final-opportunity Q snapshot in a pinned implementation and run the paired
display diagnostic. Order state, permissions and history would be identical; only the presentation as an order
record or conditional obligation would differ. For example, a buyer with value \(3\), its own standing bid of
\(1.5\) and an available ask of \(2\) gains \(1\) by crossing. With no later action available, declining to
cross gains zero.

Separate cloned contexts would test permission judgments without prompting the trading agent first. An
explicit clarification already entailed by the rules would be compared with an equal-budget generic reminder.
Display-dependent false prohibitions, costly failures to cross and correction specifically by the permission
clarification would support further investigation. Changed action rates alone would support only a framing
effect. A P comparison would also need to distinguish false prohibitions from mistaken plan sufficiency.
Complete market trials and a wider novelty check would still be needed before claiming a consequential shared
mechanism.

Separately, replaying the concentrated chip exchange and ten-point contract in P's implementation, then
inspecting the nonfinisher's final reward, would settle the smallest open question about the corrected bound.
Neither replay nor display diagnostic was completed in this thread.

\bibliographystyle{plain}
\bibliography{references}
\end{document}
